\documentclass[11pt]{letter}
\usepackage[utf8]{inputenc}
\usepackage{geometry}
\geometry{a4paper, margin=1in}
\usepackage{hyperref}

\address{Florian Rzepka\\
Technische Universität Berlin\\
Chair of Electrical Energy Storage Technology\\
Einsteinufer 11\\
10587 Berlin, Germany\\
\href{mailto:florian.rzepka@tu-berlin.de}{florian.rzepka@tu-berlin.de}}

\date{\today}

\begin{document}

\begin{letter}{Editor-in-Chief\\
Engineering Applications of Artificial Intelligence\\
Elsevier}

\opening{Dear Editor-in-Chief,}

Please find enclosed our manuscript entitled \textbf{``Embedded Artificial Intelligence in Battery Management Systems: Pruning and Quantization for Efficient State-of-Charge and State-of-Health Estimation''}, which we are submitting for exclusive consideration for publication in \textit{Engineering Applications of Artificial Intelligence}.

This work addresses the critical gap between the theoretical performance of neural network-based battery state estimators and their practical deployment on resource-constrained hardware. While many studies propose complex architectures for State-of-Charge (SOC) and State-of-Health (SOH) estimation, few rigorously quantify the computational and memory costs on actual microcontrollers.

In this study, we present a complete, reproducible pipeline that transforms Long Short-Term Memory (LSTM) models trained on the open \textit{LFP DoE Cycle Ageing Dataset} into energy- and resource-efficient AI for an STM32H7 microcontroller. We demonstrate that:
\begin{itemize}
    \item Pruning and quantization enable energy- and resource-efficient AI, reducing the Flash footprint and energy consumption by approximately 40\% with negligible loss in accuracy.
    \item Complex SOC and SOH models run autonomously on-device, as INT8 quantization halves the memory requirement to fit standard BMS hardware.
    \item The proposed estimators meet strict real-time constraints, validated through extensive hardware-in-the-loop benchmarking.
\end{itemize}

We believe this paper is highly relevant to the readership of \textit{EAAI} as it provides a concrete engineering framework for embedding AI in industrial control systems, bridging the fields of machine learning, embedded systems, and energy storage.

We confirm that this manuscript has not been published elsewhere and is not under consideration by another journal. All authors have approved the manuscript and agree with its submission.

Thank you for your consideration.

Sincerely,

Florian Rzepka\\
Technische Universität Berlin
\end{letter}
\end{document}
